\section[Cote 150 : voisinages coniques, un contre-exemple sur {[0, 1]}]{Cote 150 : voisinages coniques, un contre-exemple sur $[0, 1]$\protect\verifie{Folder150}}\label{sec:150}

La cote 150, \emph{Espaces stratifiés et voisinages côniques (ou : déploiement des espaces stratifiés)}, quatre-vingt-dix-huit pages de décembre 1981 et janvier 1982, cherche à déplier un espace stratifié : remplacer chaque strate par une version rognée, sans jointure, et garder la trace de la façon dont les morceaux se recollaient. Les pages~11 à~19 en posent le vocabulaire. Un \emph{voisinage conique} d'un fermé $Y$ dans un espace $X$ est un voisinage $T$ de $Y$ muni d'une action $(\lambda, t) \mapsto \lambda t$ du monoïde multiplicatif $\mathbf I = [0, 1]$, qui fixe $Y$ et écrase $T$ sur $Y$ pour $\lambda = 0$. Son bord $\dot T$ est la base du cône, et $f : \dot T \to Y$, $t \mapsto 0 \cdot t$, envoie chaque point du bord sur la pointe de son segment. Ôter l'intérieur de $T$ laisse un espace à bord $X^\sharp = X \smallsetminus \mathring T$, que la page note $X^{*}$, dont le bord est $\dot T$.

La page~19 énonce une proposition de reconstruction (\lot{150}{1}{19}). Les systèmes $(X, Y, (X_i)_{i \in I}, T)$, où les $X_i$ sont des fermés et $T$ un voisinage conique de $Y$ compatible avec eux, seraient équivalents, en ne gardant pour morphismes que les isomorphismes, aux systèmes $(X^\sharp, \dot X, (X_i^\sharp), Y, f : \dot X \to Y)$, où $\dot X = \dot T$ et $X_i^\sharp = X_i \cap X^\sharp$. Dans l'autre sens, on recolle à $X^\sharp$, le long de $\dot X$, le cylindre de l'application $f$. L'idée est simple : couper le voisinage conique ne perd rien, puisque le cône se reconstruit à partir de sa base et de la projection sur la pointe.

\noindent\emph{Les définitions.} Une partie $T' \subset T$ est \emph{conique} quand elle est formée des segments issus de sa trace sur le bord et de sa trace sur $Y$ : $T' = \mathbf I \cdot (T' \cap \dot T) \cup (T' \cap Y)$. C'est la remarque de la page~15 (\lot{150}{1}{15}) ; elle contient la stabilité par l'action. Le voisinage $T$ est \emph{compatible} avec un fermé $X'$ si $X' \cap T$ est conique. Un \emph{collier} de $\dot T$ dans $X^\sharp$ est un voisinage $E$ de $\dot T$ identifié à $\dot T \times [0, 1]$, et $X'$ lui est compatible si $X' \cap E$ s'écrit $X'_0 \times [0, 1]$. La partie conique $T'$ est \emph{stricte} si sa trace sur $Y$ se réduit aux pointes de ses segments, $T' \cap Y = 0 \cdot (T' \cap \dot T)$ : c'est la condition que la page~16 ajoute au bas du feuillet et entoure (\lot{150}{1}{16}).

\noindent\emph{Pourquoi la proposition demande une condition.} Les données de droite ne voient de $X_i$ que sa partie hors de l'intérieur de $T$. Dans $T$, la compatibilité reconstruit les segments à partir de leur base $X_i \cap \dot T$, mais elle laisse libre un morceau de $Y$ que $X_i$ peut contenir sans qu'aucun segment n'y aboutisse. Ce morceau est invisible à droite. L'exemple le plus petit est sur un segment.

\noindent\emph{L'exemple.} On prend $X = [0, 1]$, $Y = \{0\}$ et $T = [0, \frac12]$, avec l'action $(\lambda, t) \mapsto \lambda t$. Le bord est $\dot T = \{\frac12\}$, l'application $f$ envoie $\frac12$ sur $0$, et $X^\sharp = [\frac12, 1]$. Le collier est $E = [\frac12, \frac34]$, identifié à $\dot T \times [0, 1]$ par $(\frac12, s) \mapsto \frac12 + \frac s4$. Le voisinage $T$ est strict, au sens de la page~14 : c'est $\frac12 \cdot X$, et $X$ est lui-même un voisinage conique de $Y$ pour la même action.

\begin{encadre}
\begin{contrexemple}[contre-exemple de l'édition, page~19]
Dans l'exemple ci-dessus, les fermés $X_1 = \{0\}$ et $X_1 = \emptyset$ de $[0, 1]$ sont tous deux compatibles avec $T$ et avec le collier $E$, et ils ont la même trace, vide, sur $X^\sharp$. Les deux systèmes ont donc les mêmes données de droite, mais ils ne sont pas isomorphes, puisque l'un des fermés est vide et l'autre non. Le fermé $\{0\}$ ne vérifie pas la condition $X_1 \cap Y = f(\dot X_1)$, où $\dot X_1 = X_1 \cap \dot T$.
\end{contrexemple}
\end{encadre}

\begin{proof}
Pour $X_1 = \{0\}$ : $X_1 \cap T = \{0\}$, sa trace sur le bord est vide et sa trace sur $Y$ est $\{0\}$, de sorte que $\mathbf I \cdot \emptyset \cup \{0\} = \{0\}$ : la partie est conique. Elle ne rencontre pas $E$, et $X_1 \cap E = \emptyset = \emptyset \times [0, 1]$. Pour $X_1 = \emptyset$, tout est vide. Les deux traces sur $X^\sharp = [\frac12, 1]$ sont vides. Enfin $\{0\} \cap Y = \{0\}$, alors que $\dot X_1 = \emptyset$ et $f(\emptyset) = \emptyset$.
\end{proof}

Le contre-exemple tient encore si l'on demande que $T$ soit strict, comme le faisait la première définition de \q{strictement compatible} de la lecture, et si l'on ajoute le collier : ni l'une ni l'autre condition ne regarde la trace de $X_1$ sur $Y$. La lecture a donc mis dans la définition la condition $X_i \cap Y = f(\dot X_i)$, que la page impose elle-même au bas de la page~16. Sous cette condition, l'exemple ne présente plus d'ambiguïté.

\begin{encadre}
\begin{enonce}[énoncé de l'édition, page~19, dans l'exemple]
Dans l'exemple ci-dessus, soient $X_1$ et $X_2$ deux fermés de $[0, 1]$ compatibles avec $T$ et avec le collier, qui vérifient $X_k \cap Y = f(\dot X_k)$. Si $X_1 \cap X^\sharp = X_2 \cap X^\sharp$, alors $X_1 = X_2$.
\end{enonce}
\end{encadre}

\begin{proof}
Une partie conique $C$ de $T$ est $T$ tout entier si elle contient $\frac12$, puisqu'elle contient alors le segment $\mathbf I \cdot \frac12 = [0, \frac12]$ ; sinon sa trace sur le bord est vide et $C = C \cap Y \subset \{0\}$. Sous la condition stricte, le second cas donne $C \cap Y = f(\emptyset) = \emptyset$, donc $C = \emptyset$. Ainsi $X_k \cap T$ vaut $T$ ou $\emptyset$ selon que $\frac12$ est dans $X_k$ ou non. Or $\frac12$ appartient à $X^\sharp$, et l'on lit sur la trace $X_k \cap X^\sharp$ s'il est dans $X_k$. Comme $[0, 1] = T \cup X^\sharp$, on a $X_k = (X_k \cap T) \cup (X_k \cap X^\sharp)$, et les deux morceaux sont déterminés par la trace sur $X^\sharp$.
\end{proof}

\begin{remarque}
Le collier ne sert pas à la démonstration de l'énoncé : dans l'exemple, la condition stricte suffit. La page laisse deux réparations possibles : demander la condition stricte, ou ajouter les traces $X_i \cap Y$ aux données de droite. La lecture a retenu la première, et l'exemple montre qu'elle répare le défaut qu'il met en évidence ; il ne démontre pas la proposition corrigée en général.
\end{remarque}

\noindent\emph{Portée de la vérification.} Lean traite l'exemple seul, dans $\mathbb R$ : les parties $X$, $Y$, $T$, $\dot T$, $X^\sharp$ et le collier y sont des intervalles explicites, et les notions de partie conique, de compatibilité, de compatibilité au collier et de compatibilité stricte y sont écrites pour cette action. Il démontre que $T$ est conique et strict (\texttt{isConic\_T}, \texttt{T\_eq\_half\_smul}), le contre-exemple (\texttt{compatible\_not\_determined}, \texttt{not\_strictlyCompatible\_singleton\_zero}) et l'énoncé (\texttt{eq\_of\_strictlyCompatible}). Ne sont pas formalisés : les deux catégories de la proposition et le foncteur de recollement $X = X^\sharp \amalg_{\dot X} \mathcal C(f)$ pour des espaces quelconques, la transversalité et l'équisingularité de la même page. Mathlib n'a pas les espaces stratifiés ni les voisinages coniques.

Ne sont pas démontrés ici, de la même cote :

\begin{encadre}
\begin{theoreme*}[Grothendieck, pages 1 et 45 à 47, lissité des fibrations de drapeaux\piste{150-flag-fibration-smoothness}]\nonverifie
Dans le déploiement d'une stratification équisingulière non singulière, présenté comme foncteur sur les drapeaux de $I$, soit $d = (i_1, \dots, i_p)$ un segment initial de $d' = (i_1, \dots, i_n)$. Le bord des fibres de $\Sigma_{d'} \to \Sigma_d$ est stratifié cylindriquement par les drapeaux non vides de la somme ordinale $I_{]i_p,i_{p+1}[} \sqcup \dots \sqcup I_{]i_{n-1},i_n[}$. La fibration est donc à fibres sans bord exactement quand $i_{k+1}$ couvre $i_k$ dans $I$ pour $p \leq k < n$, condition où $i_1, \dots, i_{p-1}$ n'interviennent pas.
\end{theoreme*}
\end{encadre}
{\small Le mot qui qualifie \q{successeur} à la page~47 est une addition incertaine ; le critère est fixé par la page~46, lisible. Mathlib n'a pas les stratifications équisingulières ni leurs déploiements. On se reportera aux notes de Grothendieck (\lot{150}{1}{1}).\par}

\begin{encadre}
\begin{theoreme*}[Grothendieck, pages 47 à 66 et 70, le déploiement comme espace ordonné\piste{150-unfolding-as-ordered-space}]\nonverifie
Les données déployées d'une stratification équisingulière indexée par un ensemble ordonné fini $I$ (strates déployées $\Sigma_i$, tubes de raccord $\Sigma_{ij}$ fibrés sur $\Sigma_i$ et plongés de façon fermée, avec collier, dans $\Sigma_j$) équivalent à un seul espace topologique ordonné $\Sigma$, dont le graphe de l'ordre est la réunion des tubes. Cet espace est soumis à des conditions de propreté, de transversalité et de finitude, chaque $\Sigma_{<x}$ étant fini et totalement ordonné, et muni d'une application strictement croissante $\pi_0(\Sigma) \to I$, non nécessairement surjective.
\end{theoreme*}
\end{encadre}
{\small Les feuillets en esquissent une démonstration incomplète (pages 57 à 70). Seules la reconstitution de l'ordre et la condition~b) y sont démontrées, et l'argument de dimension des pages 67 à 69 demande des hypothèses que l'édition ajoute. Mathlib n'a pas les stratifications équisingulières. On se reportera aux notes de Grothendieck (\lot{150}{3}{47}).\par}
